\documentclass{beamer}
\usetheme{Madrid}
\usecolortheme{whale} 
\usepackage[dvipsnames]{xcolor}
\usepackage{amsmath}
\usepackage{amsfonts}
\usepackage{booktabs}
\usepackage[most]{tcolorbox}

% ---------- Color Palette ----------
\definecolor{ivorybg}{RGB}{248,246,240}
\definecolor{charcoal}{RGB}{45,45,45}
\definecolor{deepblue}{RGB}{28,70,120}
\definecolor{softteal}{RGB}{52,120,112}
\definecolor{softgold}{RGB}{191,161,111}

\setbeamercolor{normal text}{fg=charcoal,bg=ivorybg}
\setbeamercolor{structure}{fg=deepblue}
\setbeamercolor{frametitle}{fg=white,bg=deepblue}
\setbeamercolor{title}{fg=white,bg=deepblue}
\setbeamercolor{section in toc}{fg=deepblue}

\newtcolorbox{ResearchBox}[1]{
    colback=ivorybg,
    colframe=softgold,
    coltitle=white,
    colbacktitle=deepblue,
    fonttitle=\bfseries,
    title=#1,
    arc=4pt,
    boxrule=1.2pt
}

\title[QML Research]{Quantum Support Vector Machine for Prostate Cancer Detection}
\subtitle{Comparative Analysis: Classical vs. Quantum Kernel Methods}
\author{Narainkarthi}
\date{Feb 2026}

\begin{document}

\maketitle

%================================================
\section{Classical Support Vector Machine}
%================================================

%------------------------------------------------
\begin{frame}{1. Problem Setup: Linear Classification}

\begin{ResearchBox}{Fundamental Goal}
The fundamental goal of a Support Vector Machine (SVM) is to classify data points by finding an optimal boundary (Hyperplane).
\end{ResearchBox}

\begin{itemize}
    \item \textbf{Input Data:} $x \in \mathbb{R}^d$, where $d$ is the number of features.
    \item \textbf{Labels:} Binary classification using $y \in \{+1, -1\}$.
    \item \textbf{Objective:} Identify the hyperplane that separates these classes with the \textbf{maximum possible margin}.
\end{itemize}

\end{frame}

%------------------------------------------------
\begin{frame}{2. Decision Boundary and Geometry}

The hyperplane is mathematically defined by its orientation ($w$) and its position ($b$) in the feature space.

\begin{itemize}
    \item \textbf{Weight Vector ($w$):} Determines orientation and reflection of the boundary.
    \item \textbf{Bias ($b$):} Shifts the hyperplane from the origin $(0,0)$.
\end{itemize}

\begin{ResearchBox}{Fundamental Constraint}
\[
y_i(w \cdot x_i + b) \ge 1
\]
\end{ResearchBox}

\end{frame}

%------------------------------------------------
\begin{frame}{2. Decision Boundary and Geometry (Cont.)}

\textbf{Classification Regions}

\begin{itemize}
    \item \textbf{Margin Edge (SV):}
    \[
    y_i(w \cdot x_i + b) = 1
    \]
    Lies exactly on the margin boundary.
    
    \item \textbf{Positive Side:}
    \[
    y_i(w \cdot x_i + b) > 1
    \]
    Lies within the positive class region.
    
    \item \textbf{Negative Side:}
    \[
    y_i(w \cdot x_i + b) < 1
    \]
    Lies within the negative class region.
\end{itemize}

\end{frame}

%------------------------------------------------
\begin{frame}{3. Objective Function}

The goal is to maximize the margin width (distance between the two edges).

\begin{itemize}
    \item \textbf{Margin Width:}
    \[
    \frac{2}{\|w\|}
    \]
    
    \item \textbf{Optimization Objective:}
    Maximizing the margin is equivalent to minimizing the squared norm of the weights:
\end{itemize}

\[
\min \frac{1}{2} \|w\|^2 
\quad 
\text{subject to} 
\quad 
y_i(w \cdot x_i + b) \ge 1
\]

\end{frame}

%------------------------------------------------
\begin{frame}{3. Objective Function (Cont.)}

\begin{itemize}
    \item \textbf{Support Vectors (SV):}
    Only the data points lying on the margin define the hyperplane.
    
    \item The rest of the dataset does not influence the final boundary position.
    
    \item \textbf{Hard Margin:}
    Requires strict separation.
    
    \item \textbf{Soft Margin:}
    Allows some tolerance for data points inside the margin to handle noise.
\end{itemize}

\end{frame}

%================================================
\section{Kernel Methods}
%================================================

%------------------------------------------------
\begin{frame}{4. The Kernel Trick: Handling Non-Linearity}

When data cannot be separated by a straight line in its original space, we use a \textbf{Feature Map}.

\begin{itemize}
    \item \textbf{Dimension Lifting:}
    Map data to a higher-dimensional space where linear separation exists.
    
    \item \textbf{Efficiency Problem:}
    Explicit high-dimensional computation is expensive.
    
    \item \textbf{Kernel Trick:}
    Compute only the dot product (similarity) in that higher dimension.
    
    \item \textbf{Gram Matrix:}
    Similarity matrix built from kernel values.
\end{itemize}

\end{frame}

%------------------------------------------------
\begin{frame}{4. Example: Polynomial Kernel}

For 2D vectors $A$ and $B$, consider:

\[
\phi(x) = (x_1^2, x_2^2, \sqrt{2}x_1x_2)
\]

Kernel:
\[
K(A,B) = (A \cdot B)^2
\]

If:
\[
A = (1,2), \quad B = (3,4)
\]

Inner product:
\[
(1 \times 3) + (2 \times 4) = 11
\]

Kernel result:
\[
11^2 = 121
\]

\end{frame}

%================================================
\section{Quantum Perspective}
%================================================

%------------------------------------------------
\begin{frame}{5. The Balancing Act: Qubits vs Noise}

Navigating the hardware limits of current Quantum Computing.

\begin{itemize}
    \item \textbf{The Qubit Dilemma:}
    Increasing qubits increases decoherence and random noise.
    
    \item \textbf{The Gate Dilemma:}
    Deeper circuits cause cumulative error.
    
    \item \textbf{Optimization Failure:}
    Excessive noise leads to \textbf{Barren Plateaus}.
\end{itemize}

Gradients become flat and the model cannot find the optimal solution.

\end{frame}

%------------------------------------------------
\begin{frame}{6. Pre-Processing: The Gateway to Orthogonality}

Refining classical data to ensure unique and distinct quantum states.

\begin{itemize}
    \item \textbf{Standard Scaling:}
    Prevents domination of specific features.
    
    \item \textbf{PCA (Principal Component Analysis):}
    \begin{itemize}
        \item Orthogonal components
        \item Dimensionality reduction
        \item Keeps circuit shallow
    \end{itemize}
    
    \item Ensures each quantum state $|\psi(x_i)\rangle$ is distinct.
\end{itemize}

\end{frame}

%------------------------------------------------
\begin{frame}{7. Entering the Quantum World: The $[0,\pi]$ Map}

Mapping classical data to the Bloch Sphere.

\begin{itemize}
    \item $0$ = North Pole ($|0\rangle$)
    \item $\pi$ = South Pole ($|1\rangle$)
    \item Half-circle mapping ensures uniqueness
\end{itemize}

Each qubit remains independent at this stage.

\end{frame}

%------------------------------------------------
\begin{frame}{8. Transition: From Bloch Sphere to Hilbert Space}

\begin{itemize}
    \item Apply entanglement using CNOT gates
    \item Move from individual Bloch spheres to unified Hilbert space
    \item Dimensionality grows exponentially: $2^n$
\end{itemize}

Creates space for linear separation of complex non-linear data.

\end{frame}

%------------------------------------------------
\begin{frame}{9. Hybrid Strategy: Quantum + Classical}

\textbf{Quantum Strength:}

\[
K(x_i, x_j) = |\langle\psi(x_i)|\psi(x_j)\rangle|^2
\]

Measures physical overlap (similarity).

\textbf{Classical Strength:}

\begin{itemize}
    \item Receives kernel matrix
    \item Performs convex optimization
    \item Draws final hyperplane
\end{itemize}

\end{frame}

%------------------------------------------------
\begin{frame}{10. Application: Quantum Object Tracking}

\begin{itemize}
    \item 2D video frames may contain overlapping objects
    \item Classical trackers may oversimplify hidden features
    \item Hilbert space provides richer separability
\end{itemize}

\begin{ResearchBox}{Research Question}
Will quantum kernels improve real-time visual inference?
\end{ResearchBox}

\end{frame}

\end{document}